% ============================================================
%  LAMPIRAN 9 - DAFTAR RELASI BERLABEL WRONG
%  Compile standalone:  pdflatex lampiran-9.tex
%  Compile as part of thesis: pdflatex main.tex
% ============================================================
\documentclass[../../thesis]{subfiles}
\usepackage{hyphenat}
\usepackage{iftex}

% -- Encoding, Language & Fonts --------------------------------
\ifXeTeX
    \usepackage{polyglossia}
    \setmainlanguage{bahasai}
    \usepackage{fontspec}
    \setmainfont{TeX Gyre Termes}
    \setsansfont{TeX Gyre Heros}
    \setmonofont[Scale=0.88]{TeX Gyre Cursor}
\else
    \usepackage[utf8]{inputenc}
    \usepackage[T1]{fontenc}
    \usepackage[indonesian]{babel}
    \usepackage{mathptmx}
    \usepackage{helvet}
    \usepackage{courier}
\fi

% -- Page geometry (A4, UI thesis margins) ---------------------
\usepackage[
  a4paper,
  headheight = 0pt,
  top    = 3cm,
  bottom = 3cm,
  left   = 3cm,
  right  = 3cm
]{geometry}

% -- Line spacing & paragraph ----------------------------------
\usepackage{setspace}
\onehalfspacing
\setlength{\parindent}{1.27cm}
\setlength{\parskip}{0pt}

% -- Microtype & justification ---------------------------------
\usepackage{microtype}
\sloppy

% -- Headings ---------------------------------------------------
\usepackage{titlesec}

\titleformat{\chapter}[block]
  {\centering\rmfamily\bfseries\large}{}
  {0pt}{\MakeUppercase}
\titlespacing*{\chapter}{0pt}{0pt}{1.5em}

\titleformat{\section}
  {\rmfamily\bfseries\normalsize}{\thesection}{1em}{}
\titlespacing*{\section}{0pt}{1.2em}{0.6em}

\titleformat{\subsection}
  {\rmfamily\bfseries\normalsize}{\thesubsection}{1em}{}
\titlespacing*{\subsection}{0pt}{1em}{0.5em}

% -- Tables -----------------------------------------------------
\usepackage{longtable, booktabs, array}
\usepackage{tabularx}

% -- Hyperlinks -------------------------------------------------
\usepackage[hidelinks]{hyperref}

% ============================================================
\begin{document}

% When standalone: switch to appendix mode and set the counter
% so appendix-local numbering renders correctly. When compiled inside main.tex,
% main.tex issues \appendix and advances the counter before \subfile.
\ifSubfilesClassLoaded{}{\appendix\setcounter{chapter}{9}}

\lampiran{Daftar Relasi Berlabel \texttt{wrong}}
\label{lamp:wrong-triples}

Lampiran ini memuat seluruh relasi semantik intra-mapel yang memperoleh
konsensus akhir \texttt{wrong} pada validasi pakar Fase~1. Relasi-relasi ini
tetap dipertahankan dalam graf sebagai artefak audit dan analisis galat, tetapi
dapat dikeluarkan melalui properti \texttt{consensus}/\texttt{status} untuk uji
ketahanan metrik struktural pada Tabel~\ref{tab:struktural-robustness}. Daftar
berikut mencakup seluruh 21 relasi berlabel \texttt{wrong}; uji ketahanan
struktural hanya memengaruhi relasi yang berada pada subgraf
\texttt{Concept}$\rightarrow$\texttt{Concept}.

\begingroup
\scriptsize
\setlength{\tabcolsep}{3pt}
\begin{longtable}{@{}
  >{\centering\arraybackslash}p{0.05\linewidth}
  >{\raggedright\arraybackslash}p{0.10\linewidth}
  >{\raggedright\arraybackslash}p{0.22\linewidth}
  >{\raggedright\arraybackslash}p{0.57\linewidth}@{}}
\caption{Relasi semantik intra-mapel dengan konsensus akhir \texttt{wrong}.}
\label{tab:wrong-triples}\\
\toprule
\textbf{No.} & \textbf{Mapel} & \textbf{Bab} & \textbf{Relasi} \\
\midrule
\endfirsthead
\toprule
\textbf{No.} & \textbf{Mapel} & \textbf{Bab} & \textbf{Relasi} \\
\midrule
\endhead
1 & Biologi & Evolusi &
Seleksi Alam -- \texttt{TERDIRI\_DARI} -- Seleksi Alami \\
2 & Fisika & KEMAGNETAN &
Medan Magnet Induksi Kawat Melingkar Berarus -- \texttt{MENGATUR} --
Arah Medan Magnet (Kawat Melingkar) \\
3 & Fisika & RELATIVITAS &
Transformasi Lorentz -- \texttt{MENGATUR} -- Faktor Lorentz ($\gamma$) \\
4 & Kimia & ELEKTROKIMIA &
Bilangan Oksidasi -- \texttt{MEMPENGARUHI} -- Perubahan Kimia \\
5 & Kimia & ELEKTROKIMIA &
Oksidasi -- \texttt{MENGHASILKAN} -- Agen Pengoksidasi \\
6 & Kimia & ELEKTROKIMIA &
Pelapisan Logam (Electroplating) -- \texttt{MEMUNGKINKAN} --
Pemurnian Tembaga \\
7 & Kimia & ELEKTROKIMIA &
Reduksi -- \texttt{MENGHASILKAN} -- Agen Pereduksi \\
8 & Kimia & GUGUS FUNGSI DALAM SENYAWA KARBON &
Reaksi Adisi -- \texttt{BEREAKSI\_DENGAN} -- Alkena \\
9 & Kimia & GUGUS FUNGSI DALAM SENYAWA KARBON &
Reaksi Esterifikasi -- \texttt{BEREAKSI\_DENGAN} -- Asam Karboksilat \\
10 & Kimia & GUGUS FUNGSI DALAM SENYAWA KARBON &
Reaksi Oksidasi Senyawa Organik -- \texttt{DIKATALISIS\_OLEH} --
Piridinium Klorokromat (PCC) \\
11 & Kimia & GUGUS FUNGSI DALAM SENYAWA KARBON &
Reaksi Reduksi Senyawa Organik -- \texttt{DIKATALISIS\_OLEH} --
Litium Alumunium Hidrida (LiAlH4) \\
12 & Kimia & GUGUS FUNGSI DALAM SENYAWA KARBON &
Reaksi Substitusi -- \texttt{BEREAKSI\_DENGAN} -- Alkil Halida \\
13 & Kimia & LARUTAN DAN KOLOID &
Asam Lemah -- \texttt{DIFORMULASIKAN\_SEBAGAI} -- Ka \\
14 & Kimia & LARUTAN DAN KOLOID &
Asam Lemah -- \texttt{DIFORMULASIKAN\_SEBAGAI} -- $\alpha$ (Derajat
Ionisasi) \\
15 & Kimia & LARUTAN DAN KOLOID &
Basa Lemah -- \texttt{DIFORMULASIKAN\_SEBAGAI} -- $\alpha$ (Derajat
Ionisasi) \\
16 & Kimia & LARUTAN DAN KOLOID &
Basa Lemah -- \texttt{DIFORMULASIKAN\_SEBAGAI} -- Kb \\
17 & Kimia & LARUTAN DAN KOLOID &
pH -- \texttt{MENGATUR} -- Keasaman Larutan \\
18 & Kimia & MAKROMOLEKUL ORGANIK &
Degradasi Plastik -- \texttt{TERDIRI\_DARI} -- Biodegradable Plastic \\
19 & Kimia & MAKROMOLEKUL ORGANIK &
Degradasi Plastik -- \texttt{TERDIRI\_DARI} -- Photodegradable Plastic \\
20 & Kimia & MAKROMOLEKUL ORGANIK &
Gaya Antarmolekul Polimer -- \texttt{MENYEBABKAN} -- Polimer Termoset \\
21 & Kimia & MAKROMOLEKUL ORGANIK &
Mekanisme Polimerisasi Adisi -- \texttt{MENYEBABKAN} -- Polimerisasi
Adisi \\
\bottomrule
\end{longtable}
\endgroup

\end{document}
